\documentclass[exercice]{thedocument}%
\usepackage{texmaths-tikz}%
\startdatakeys
\source{}
\cycle{}
\competencesDuSocle{}
\connaissancesRequises{}
\competencesTravaillees{}
\enddatakeys
\begin{document}%
\begin{enumerate}
    \item Voici la représentation graphique d'une fonction $f$ :\par\noindent
    \begin{tikzpicture}[baseline=(O.base),scale=0.45]
        \coordinate(O)at(0,5);
        \simpleFunc xmin=-5,xmax=5,ymin=-2,ymax=9,gridstep=1,func=f=x^2/4;
    \end{tikzpicture}\par\noindent
    Par lecture graphique :
    \begin{enumerate}
        \item donner la valeur de $f(2)$.
        \item quelle est l'images de -4 par la fonction $f$~?
        \item donner un maximum d'antécédents de 4 par $f$.
    \end{enumerate}
    \item L'expression algébrique de cette fonction est\par\noindent
    {\centering$f(x)=\dfrac{x^2}4$.\par}\noindent
    Calculer $f(3)$ puis déterminer les antécédents de~5~par~$f$.
\end{enumerate}
\end{document}
